% answers_short.tex —— 1992 年·三南卷 · 简短答案速查
% 由 tools/make_short_answers.py 自 answers.tex 的 \gksolhead/\gkanswer 提取，
% 并与 1992/1992_数学真题_三南_答案要点.md 逐题比对后生成，勿手改（改源 MD 后重跑）。

\begin{gknotice}
\textbf{使用说明}：本节只给最终答案便于快速核对，完整推导见解析版。答案为编写者独立推导，\textbf{非官方标准答案}。
\end{gknotice}

\gkpart{选择题}
\begin{center}\begin{tabular}{cccccccccccccccccc}
\toprule
题号 & 1 & 2 & 3 & 4 & 5 & 6 & 7 & 8 & 9 & 10 & 11 & 12 & 13 & 14 & 15 & 16 & 17 \\
\midrule
答案 & C & D & A & B & C & A & D & B & C & A & A & D & B & B & D & C & A \\
\bottomrule
\end{tabular}\end{center}

\gkpart{填空题}
\begin{itemize}[leftmargin=2.6em]
  \item[\textbf{18.}] \(\sqrt2-1\)
  \item[\textbf{19.}] \(y=\sqrt3x+3-2\sqrt3\)
  \item[\textbf{20.}] \((x+3)^2+(y-3)^2=9\)
  \item[\textbf{21.}] \(\dfrac13\)
  \item[\textbf{22.}] \(81\)
  \item[\textbf{23.}] \(\dfrac{3aV}{2S_1S_2}\)
\end{itemize}

\gkpart{解答题}
\begin{itemize}[leftmargin=2.6em]
  \item[\textbf{24.}] \(a=\dfrac12\) 时另有根 \(x=1-\log_23\)；\(a=3\) 时另有根 \(x=1-\log_32\)。
  \item[\textbf{25.}] 见解析（线面平行判定）。
  \item[\textbf{26.}] 见解析（逐项放缩后望远镜求和）。
  \item[\textbf{27.}] \((y-2)^2=32\left(x+\dfrac32\right)\)，即 \(y^2-4y-32x-44=0\)。
  \item[\textbf{28.}] \(z=1\pm3\mathrm{i}\) 或 \(z=3\pm\mathrm{i}\)。
\end{itemize}
